\section{测试策略：端到端验证}

\subsection{过早标记完成的问题}

Anthropic 团队观察到的另一个主要失败模式是：Agent 倾向于在没有充分测试的情况下将功能标记为完成。典型的行为是——Agent 修改了代码，可能还跑了单元测试或发了几个 curl 请求，但没有验证功能是否在端到端场景下真正可用。

\subsection{浏览器自动化测试}

在 Web 应用开发场景下，Anthropic 团队发现让 Agent 使用浏览器自动化工具（如 Puppeteer MCP）进行端到端测试可以显著提升质量。Agent 被要求"像人类用户那样"测试功能：

\begin{itemize}
    \item 在浏览器中打开应用
    \item 执行用户操作（点击、输入、导航）
    \item 验证 UI 响应是否符合预期
    \item 检查状态变更是否正确持久化
\end{itemize}

\begin{importantbox}{测试工具决定测试质量}
提供合适的测试工具对于 Agent 的表现至关重要。当 Agent 只能通过代码审查和单元测试来验证时，它很难发现那些只在端到端场景下才暴露的 bug。而当给予 Agent 浏览器自动化能力后，它能够像真实用户一样操作应用，发现并修复仅通过代码分析无法识别的问题。\textbf{工具的能力边界决定了 Agent 的能力边界}。
\end{importantbox}

\subsection{测试工具的局限性}

尽管浏览器自动化测试大幅提升了 Agent 的验证能力，但仍然存在一些限制：

\begin{itemize}
    \item Claude 的视觉能力有其局限——某些 UI 细节可能无法被准确识别
    \item 浏览器原生组件（如 \texttt{alert} 弹窗）在 Puppeteer MCP 中不可见
    \item 依赖这类不可见组件的功能往往表现出更多 bug
    \item 复杂的动画、过渡效果难以通过自动化测试验证
\end{itemize}

\begin{knowledgebox}{工具链的可观测性}
这个发现揭示了 Agent 系统设计中一个重要原则：Agent 的调试和测试能力受限于其工具链的可观测性（observability）。如果某些系统行为对 Agent 不可见（如原生弹窗、后台定时任务、网络延迟等），Agent 就无法有效地测试和修复相关问题。在设计 Agent harness 时，应该尽量扩展 Agent 的可观测范围，或者在 prompt 中明确告知这些盲区。
\end{knowledgebox}

\subsection{本章小结}

端到端测试是确保 Agent 不过早标记功能完成的关键。通过浏览器自动化工具，Agent 可以像真实用户一样测试 Web 应用。但工具链的局限性意味着某些类型的 bug 仍然难以被自动发现，这是未来需要继续改进的方向。


